\documentclass[a4paper, 12pt]{extarticle}
\usepackage[utf8]{inputenc}
\usepackage{amsmath,amssymb, amsthm}
\usepackage{enumitem}
\usepackage{listings}
\lstset{basicstyle=\ttfamily}

\newtheorem*{statement}{Statement}

\begin{document}
\lstMakeShortInline`

\begin{enumerate}
\item `x == 1, r_y == 0, r_z == 0` \\
$A.1 \to B.1 \to A.2 \to B.2 \to A.3 \to B.3 \to A.4 \to B.4 \to A.5 \to B.5$.
on `A.5` and `B.5` a data race occurs! In this case, we write the same value(1) to the `x`, just twice, therefore `x` will be 1. `y` and `z` remain 0, so `r_y` and `r_z` are zeroed as well.

\item `x == 2, r_y == 0, r_z == 1` \\
$A.1 \to B.1 \to B.2 \to A.2 \to A.3 \to B.3 \to A.4 \to A.5 \to B.4 \to B.5$.
Using the fact that `x = 0`, we pass both `if`'s, then, in thread B, we read `x` into `z` only after updating `x` in the thread A, getting `z = 1; x = 1 + 1; r_z = 1`.
\item `x == 1, r_y == 0, r_z == 1` \\
\begin{statement}
The trace of `x == 1, r_y == 0, r_z == 1` is impossible to occur in our case.
\end{statement}
\begin{proof}
Let's assume there exists concurrent execution trace where `main` thread observes `x == 1, r_y == 0, r_z == 1`.
\begin{enumerate}
\item Then $B.5$ happened earlier than $A.5$. (for `r_y = y == 0`, `x` at $A.4$ must be equal to 0, therefore $A.5$ will write 1 to `x`. only after getting `r_y = 0` we can go to $B.5$, where we will get `x = 2`! However, we can pass that by executing $A.5$ after the $B.5$, so `x` becomes 1.)
\item Then $A.5$ happened earlier than $B.4$. (we need to have `x` set to 1 to set `z` to 1 to set `r_z` to 1.).
\item But also, according to program order, $B.4$ always happens earlier than $B.5$.
\end{enumerate}
Contradiction: $B.5$ happened earlier than $B.4$ vs. $B.4$ happened earlier than $B.5$. \\
Therefore, it is impossible to observe `x == 1, r_y == 0, r_z == 1` when using interleaving model of execution.
\end{proof}

\end{enumerate}

\end{document}
